% Reference entries, A (PUB-REF-0001 AN-11; model: WordPerfect 5.1
% Reference, feature entries pp. 31-711, docs/pubs/originals/wordperfect-
% 51.md s4).  Law 1: os/iword/app/face_a.c M_FONT / M_APPEAR (Ctrl-F8, 2,
% the nine choices; spec/iword/reveal_codes.h IW_APPEAR_MENU,
% IW_APPEAR_ATTRS), os/iword/core/codes.c iw_ed_attr / iw_attr_faces (the
% pair of codes, the cursor between them, the faces drawn).
% GATES: test-pubs-iword test-iword-codes test-iword-reveal-codes test-flair-iword-codes
\Gates{test-pubs-iword test-iword-codes test-iword-reveal-codes test-flair-iword-codes}
\begin{Feature}{Appearance}
Appearance changes the look of text without changing its size: bold,
underline, double underline, italic, outline, shadow, small capitals,
redline and strikeout.  You choose an appearance from the Font menu.

InitechWord puts an appearance in the document as a pair of codes: an
\emph{on} code where the appearance begins and an \emph{off} code where it
ends, for example \Code{[ITALC]} and \Code{[italc]}.  The text between the
two codes has the appearance.  Bold and underline have function keys of
their own (see \InRef{Bold} and \InRef{Underline}); the other appearances
are chosen only from this menu.

To give text an appearance as you type it,
\begin{WPSteps}
\item Press \FK{Font}{Ctrl-F8}, then select \Ch{Appearance}{2}.  The
Appearance menu appears on the status line:
\StatusLine{1 Bold 2 Undln 3 Dbl Und 4 Italc 5 Outln 6 Shadw 7 Sm Cap 8 Redln 9 Stkout: 0}
\item Select the appearance by its number.
\item Type the text.
\item Press \KC{Ctrl-F8}{2} and select the same appearance again to move the
cursor past its off code, so that what you type next is ordinary text.
\Or
Press \K{Right} to move past the off code.
\end{WPSteps}
\begin{WPNote}To give an appearance to text you have already typed, block
the text first (see \InRef{Block, Attributes}).\end{WPNote}

\Documents{wpmenu}{Ctrl-F8,2=Appearance}
\Documents{wpmenu}{Ctrl-F8,2,1=Bold}
\Documents{wpmenu}{Ctrl-F8,2,2=Undln}
\Documents{wpmenu}{Ctrl-F8,2,3=Dbl Und}
\Documents{wpmenu}{Ctrl-F8,2,4=Italc}
\Documents{wpmenu}{Ctrl-F8,2,5=Outln}
\Documents{wpmenu}{Ctrl-F8,2,6=Shadw}
\Documents{wpmenu}{Ctrl-F8,2,7=Sm Cap}
\Documents{wpmenu}{Ctrl-F8,2,8=Redln}
\Documents{wpmenu}{Ctrl-F8,2,9=Stkout}
\par\addvspace{6pt}
The appearances and their codes are:
\par\addvspace{4pt}
\noindent\begin{tabular}{@{}lll@{}}
\textbf{Selection} & \textbf{Appearance} & \textbf{Codes}\\
Bold (1) & bold & \Code{[BOLD]} \Code{[bold]}\\
Undln (2) & underline & \Code{[UND]} \Code{[und]}\\
Dbl Und (3) & double underline & \Code{[DBL UND]} \Code{[dbl und]}\\
Italc (4) & italic & \Code{[ITALC]} \Code{[italc]}\\
Outln (5) & outline & \Code{[OUTLN]} \Code{[outln]}\\
Shadw (6) & shadow & \Code{[SHADW]} \Code{[shadw]}\\
Sm Cap (7) & small capitals & \Code{[SM CAP]} \Code{[sm cap]}\\
Redln (8) & redline & \Code{[REDLN]} \Code{[redln]}\\
Stkout (9) & strikeout & \Code{[STKOUT]} \Code{[stkout]}\\
\end{tabular}
\par\addvspace{4pt}

For example, to type a sentence with one word in italic,
\BeginReplay{ref-appearance}
\begin{WPSteps}
\item Type \TypeRun{Read the }(with the space after it).
\item Press \Press{Font}{Ctrl-F8}, select \Choose{Appearance}{2}, then
select \Choose{Italc}{4}.
\item Type \TypeRun{new}.
\item Press \Press{Font}{Ctrl-F8}, select \Choose{Appearance}{2}, and select
\Choose{Italc}{4} again: the cursor moves past \Code{[italc]}.
\item Type \TypeRun{ policy.}(with the space before it).
\item Press \Press{Reveal Codes}{Alt-F3} to see the codes.
\end{WPSteps}
\WordScreen[rows=13]{ref-appearance}{Italic text on the editing screen and in Reveal Codes}
\EndReplay

\Notes
\Topic{Appearance on the Screen}
On the editing screen, bold, underlined and italic text is drawn in that
appearance, so you see it as you type.  Text with any other appearance --
double underline, outline, shadow, small capitals, redline or strikeout --
is drawn as ordinary text.  The appearance is in the document all the same;
press Reveal Codes to see its codes.

\Topic{Appearances Together}
Text can have several appearances at once.  Choose a second appearance
while the cursor is between the codes of the first: the second pair of codes
is put inside the first, as in \Code{[BOLD][ITALC]urgent[italc][bold]}.

\Topic{Choosing an Appearance That Is On}
If the cursor is already in text that has the appearance you select,
InitechWord does not start it again: it moves that appearance's off code to
the cursor, so the appearance ends there, and places the cursor after it.
When the off code is right at the cursor, this simply moves the cursor past
it.

\Topic{Deleting an Appearance}
To take an appearance away, delete either of its codes.  In Reveal Codes,
Backspace or Delete on the code deletes it; on the editing screen
InitechWord first asks \Scr{Delete [ITALC]? No (Yes)} (with the code's own
name), and Yes deletes it.  Either way the other code of the pair goes with
it.  See \InRef{Codes}.

\Topic{Normal}
Select \Ch{Normal}{3} from the Font menu to end every appearance and size in
effect at the cursor at once (see \InRef{Font}).

\SeeAlso{Block, Attributes; Bold; Codes; Font; Italics; Reveal Codes; Size; Underline}
\end{Feature}
